\subsection*{Функции потерь}
Для начала надо понимать, а как мы хотим измерять отклонение предсказаний
нашей обученной модели от реальных ответов? Ранее мы это определяли, как
сумма расстояний каждой конкретной точки (объекта) до построенной прямой
(или гиперплоскости). Осталось дело за малым "--- формализовать 
наши выводы. С точки зрения геометрии, конечно, это и звучит легко, 
однако на практике встречаются разные подходы к оцениванию ошибок.
Почему это так и почему не хватает какого"=то одного подхода "---
мы рассмотрим далее.

Итак, под \textit{функцией потерь} (или функционалом потерь) мы понимаем
некоторую функцию, которая принимает на вход модель, реальные ответы, и 
на этой основе считает какую-то \textit{ошибку}, т.е. это самое 
отклонение. Как уже и упоминалось, таких функций далеко не одна, и даже 
не 10, поэтому постепенно будем проходиться по наиболее популярным и 
используемым.

\subsubsection*{\texttt{MSE}, \textt{Mean Squared Error}}

Первая на очереди "--- средняя квадратичная ошибка (или \texttt{MSE},
или \texttt{Mean Squared Error}). Интуиция: давайте будем считать разницу
между реальным ответом и предсказанным, и возводить её в квадрат.

Тогда ошибка на одном объекте считается так:
\[
Q(a, y) = (a(x) - y)^2
\] 
т.е. выход модели на этом объекте минус реальный ответ на этом объекте.

Тогда для всей выборки:
\[
MSE = \frac{1}{\ell}\sum_{i=1}^\ell (a(x_i) - y_i)^2
\]

Интуиция заключается в том, чтобы измерять среднюю величину ошибок
модели. То есть мы считаем сумму всех таких ошибок, и делим их на 
количество объектов.

Ну и вот, казалось бы, мы написали прекрасный функционал потерь, почему
бы просто не использовать его везде? Давайте порассуждаем в контексте 
задачи предсказания цены квартиры: мы вычитаем из предсказанной стоимости
реальную, возводим в квадрат и получаем... рубли в квадрате? Я думаю
мало кто сможет объяснить, что такое рубли в квадрате, и своему заказчику
тем более.

\subsubsection*{\texttt{RMSE}, \texttt{Root Mean Squared Error}}

Интуиция простая донельзя "--- давайте просто возьмём корень из 
\texttt{MSE}:
\[
RMSE = \sqrt{MSE}
\]

Вот мы и решили ранее описанную проблему "--- ошибка теперь в рублях, 
теперь то точно всё хорошо... ведь всё хорошо?

Дело в том, что \texttt{MSE} прекрасно подходит для сравнения 
двух разных построенных моделей, однако по этому значению тяжело сказать,
насколько хороша одна конкретная модель. Давайте предположим, что 
в наших данных целевая переменная $y \in [0, 1]$. Если 
\texttt{MSE} выдаёт значение $10$, можно догадаться, что это огромная ошибка. При этом если целевая переменная $y \in [10000, 100000]$, то 
то же значение \texttt{MSE} выглядит очень даже неплохо.

То есть, \texttt{MSE} измеряется в квадратах единиц исходной переменной, поэтому она чувствительна к масштабу. Это означает, 
что при изменении диапазона целевой переменной та же ошибка будет оцениваться по"=разному, что может затруднить интерпретацию
в некоторых моментах.

\subsubsection*{Коэффициент детерминации}
Решая вышепоставленную проблему, напрашивается решение: давайте как"=нибудь сделаем так, чтобы \texttt{MSE} лежал в интервале 
$[0, 1]$, вне зависимости от масштаба целевой переменной. 
Так появился \textit{коэффициент детерминации}, который выглядит довольно
страшно:
\[
R^2 = 1 - \frac{\sum_{i=1}^\ell (a(x_i) - y_i)^2}{\sum_{i=1}^\ell (y_i - \overline{y})^2}
\]
На самом деле, в числителе и в знаменателе 
должен быть коэффициент $\frac{1}{\ell}$, но именно поэтому он и сократился.
Тогда если приглядеться, то в числителе "--- уже знакомый \texttt{MSE}, а в знаменателе пока какое"=то нечто. Тут $\overline{y}$ обозначает среднее значение целевой переменной. Просто представьте, что в знаменателе тоже 
\texttt{MSE} модели, которая при любом раскладе выдаёт среднее значение 
целевой переменной. При такой математике получается нормализованное (т.е. 
в диапазоне $[0, 1]$) значение \texttt{MSE}. Чем ближе результат
коэффициента детерминации к $1$, тем считается лучше. Иногда в 
литературе коэффициент детерминации называют \textit{долей объяснённой дисперсии}.

\subsection*{Лирическое отступление}
Сейчас мы игнорировали одну большую проблему "--- понятие выбросов в 
данных. Мы проиллюстрируем это подробнее чуть позже, но проблема заключается в том, что наша модель выдаст какое"=то предсказание 
на выбросе, разница возведётся в квадрат и получится гигантское значение.
По тому же коэффициенту детерминации, например, очень сложно понять,
что выброс нам что"=то испортил. Если забегать сильно вперёд, то 
по значениям этих функционалов в дальнейшем будут обновляться веса
модели "--- именно поэтому это является сильной проблемой. Наша модель
может сильно испортиться даже из за одного выброса. Эту проблему 
надо как"=то проблему.

\subsubsection*{\texttt{MAE}, \texttt{Mean Absolute Error}}
Ранее описанная проблема характеризуется как \textit{неустойчивость к
выбросам} в данных. \texttt{MAE} является наиболее устойчивым 
функционалом потерь, в чём мы убедимся позже. Давайте посмотрим на 
формулу:
\[
MAE = \frac{1}{\ell}\sum_{i=1}^\ell |a(x_i) - y_i|
\]
Здесь уже больше устойчивости к выбросам. Самое время проиллюстрировать
это на примере:
\begin{table}[H]
\begin{tabular}{|c|c|c|l|}
\hline
Ответ & Предсказание & MAE & MSE  \\ \hline
1     & 2            & 1   & 1    \\ \hline
2     & 3            & 1   & 1    \\ \hline
100   & 3            & 97  & 9409 \\ \hline
4     & 4            & 0   & 0    \\ \hline
\end{tabular}
\end{table}

\[
MSE \approx 3137
\]
\[
MAE \approx 33
\]

\texttt{MAE} отдаёт меньший приоритет выбросам, что делает эту метрику 
наиболее стабильной.

Давайте теперь подумаем: что если нам не столь важно предсказать прям 
значение ответа, а скажем, просто его порядок? Например в задаче предсказания расстояния между галактиками "--- будет не сильно страшно в 
текущих реалиях, если мы ошибёмся на пару сотен световых лет.

\subsubsection*{\texttt{MSLE}, \texttt{Mean Squared Logarithmic Error}}
Когда речь заходит о порядке какой"=то величины, обычно в этом же 
контексте упоминаются логарифмы. Действительно, давайте просто навесим
логарифм:
\[
MSLE = \frac{1}{\ell}\sum_{i=1}^\ell (\log(a(x_i) + 1) - \log(y+1))^2
\]
Надеюсь не трудно догадаться, что такое \texttt{RMSLE}:
\[
RMSLE = \sqrt{MSLE}
\]

Давайте пойдём ещё дальше. Теперь мы хотим получать не просто какое"=то
число, которое интерпретируется как какая"=то там ошибка, а хотим получать какую"=то относительную ошибку. Например сказать <<ошибка
50\%>> звучит лучше, чем <<\texttt{MSE} $= 123.41286...$>>.

\subsubsection*{\texttt{(S)MAPE}, \texttt{(Symmetric) Mean Absolute Percentage Error}}

Формула:
\[
MAPE = \frac{1}{\ell}\sum_{i=1}^\ell \frac{|y_i - a(x_i)|}{y_i} * 100\%
\]

Однако у \texttt{MAPE} есть проблема, связанная с несимметричностью 
данной функции. Давайте допустим, что ответ на объекте равен $1$,
а прогнозы нашей модели неотрицательные. При \textit{недопрогнозе}
(когда модель выдала значение меньше, чем ответ) максимальная ошибка может быть 100\%:
\[
\frac{|1 - 0|}{1} * 100\%= 100\%
\]
В случае \textit{перепрогноза} (когда модель выдала значение больше, 
чем ответ), максимальная ошибка неограничена:
\[
\frac{|1 - 10|}{1} * 100\% = 900\%
\]
Эту проблему решает \texttt{SMAPE}:
\[
SMAPE = \frac{1}{\ell}\sum_{i=1}^\ell \frac{|y_i - a(x_i)|}{(|y_i| + |a(x_i)|)/2}
\]

Функцоналов потерь, конечно, существует достаточно много, но мы 
остановимся лишь на этих базовых. Главное, что нужно из всего этого
извлечь "--- \textbf{с помощью функционала ошибок мы задаём требования
к нашей модели}.

Вернёмся к задаче обучения модели. В процессе подбора параметров 
и подсчёта ошибок, понятное дело мы хотим, чтобы выбранный функционал 
был как можно меньше. Формализуя (на примере \texttt{MSE}):
\[
\frac{1}{\ell}\sum_{i=1}^\ell(a(x_i) - y_i))^2 \rightarrow \min_w 
\]
Формально, это и есть процесс обучения модели. Мы пока не будем затрагивать процесс подбора этих параметров $w$, потому что это
сложная тема для отдельной лекции.